% V4 J1 round fixture (ADR 0082 §4): an acmart sigconf submission as a real
% one is shaped -- `review` line numbers, the conference running head, and the
% references opened on a fresh page by \clearpage. Its PDF is what
% `computePdfFacts` must read as "references begin on page 3, main text 2".
\documentclass[sigconf, screen, review, anonymous]{acmart}

\acmConference[FIXTURE '26]{The Fixture Conference}{November 2026}{Springfield, USA}
\setcopyright{none}

\newcommand{\para}{Long-context sequence models spend most of their compute on
attention, and the share grows with the length of the input. A routing scheme
that bounds the number of keys each query may read changes that growth from
quadratic to linear, at the price of a scorer that has to decide which keys are
worth reading. This paragraph stands in for the prose of a real submission: it
is long enough to fill a column with ordinary text, and it says nothing a
reviewer of layout would need to read closely.}

\title{Corridor Attention: A Fixture for the References Page}
\author{A. Fixture}
\affiliation{%
  \institution{Fixture Institute of Computing}
  \city{Springfield}
  \country{USA}}

\begin{document}

\begin{abstract}
A fixture whose only job is its page layout: two pages of body text, and a
reference list that opens the third page under the running head.
\end{abstract}

\keywords{fixtures, page layout}

\maketitle

\section{Introduction}
\para{} Prior work bounds the interaction pattern~\cite{child2019sparse}.

\para

\para

\section{Method}
\para

\para

\para

\para

\section{Experiments}
\para

\para

\para

\para

\section{Conclusion}
\para{} The conclusion ends here, on the second page.

\clearpage
\bibliographystyle{ACM-Reference-Format}
\bibliography{refs}

\end{document}
